\section{Método}\label{sec:met}

Para la realización de este estudio, la metodología se ha estructurado de la siguiente manera:

\subsection{Formulación Matemática del Problema}\label{subsec:formulacion}

Sea $G = (V, E)$ un grafo no dirigido con conjunto de vértices $V = \{1, 2, \ldots, n\}$ y conjunto de aristas $E$. Se define $\bar{E} = \{(i,j) : i,j \in V,\ i < j,\ (i,j) \notin E\}$ como el conjunto de pares no adyacentes (cada par considerado una sola vez), es decir, los pares de vértices que \textbf{no} están conectados en $G$.

\subsubsection{Formulación como Programa Lineal Entero (ILP)}

El MCP se formula como un problema de Programación Lineal Entera Binaria (Binary Integer Linear Program, BILP). Se introduce una variable de decisión binaria $x_i$ para cada vértice $i \in V$, donde $x_i = 1$ si el vértice $i$ pertenece al clique, y $x_i = 0$ en caso contrario \parencite{pardalos1994maximum, nemhauser1988}.

\vspace{0.3cm}
\noindent\textbf{Función objetivo:}
\begin{equation}\label{eq:obj}
    \max \quad z = \sum_{i \in V} x_i
\end{equation}

\noindent\textbf{Restricciones de adyacencia:}
\begin{equation}\label{eq:adj}
    x_i + x_j \leq 1, \quad \forall\ (i,j) \in \bar{E}
\end{equation}

\noindent\textbf{Restricciones de integralidad:}
\begin{equation}\label{eq:bin}
    x_i \in \{0, 1\}, \quad \forall\ i \in V
\end{equation}

La función objetivo \eqref{eq:obj} maximiza el número de vértices seleccionados. Las restricciones \eqref{eq:adj} garantizan que si dos vértices $i$ y $j$ \textbf{no} son adyacentes (es decir, $(i,j) \in \bar{E}$), entonces ambos no pueden pertenecer simultáneamente al clique. Las restricciones \eqref{eq:bin} aseguran que las variables sean binarias. Toda solución factible de este modelo corresponde a un clique en $G$, y la solución óptima es el clique de cardinalidad máxima \parencite{pardalos1994maximum, bomze1999}.

\subsubsection{Relajación Lineal Continua (LP)}

Para obtener cotas superiores del valor óptimo entero, se resuelve la relajación lineal del problema \eqref{eq:obj}--\eqref{eq:bin}, reemplazando la restricción de integralidad por su envolvente continua \parencite{nemhauser1988, schrijver1986}:

\begin{equation}\label{eq:relax}
    0 \leq x_i \leq 1, \quad \forall\ i \in V
\end{equation}

Sea $z^*_{LP}$ el valor óptimo de la relajación y $z^*_{ILP}$ el del problema entero original. Dado que la región factible del LP contiene a la del ILP, se cumple la relación fundamental:
\begin{equation}\label{eq:cota}
    z^*_{ILP} \leq z^*_{LP}
\end{equation}

Esta propiedad es la base teórica del método de Ramificación y Acotamiento: si la cota superior $z^*_{LP}$ de un subproblema no supera la mejor solución entera encontrada hasta el momento, dicho subproblema puede descartarse sin pérdida de optimalidad (\textit{poda por cota}) \parencite{nemhauser1988, wolsey1998}.

\subsection{Procedimiento de Resolución}

\begin{enumerate}
    \item \textbf{Herramientas de desarrollo:} Los experimentos se implementaron en Python 3.12, utilizando Pyomo \parencite{pyomo2021} como framework de modelado algebraico y HiGHS \parencite{highs2024} como solver MIP. Python ofrece una robusta capacidad para la formulación rápida de modelos de optimización y el manejo eficiente de matrices de adyacencia.

    \item \textbf{Instancia de prueba:} Se utilizó la instancia \textit{keller4} del 2nd DIMACS Implementation Challenge \parencite{dimacs1993, mascia_dimacs}, un grafo de 171 vértices y 9\,435 aristas ampliamente reconocido en la literatura para evaluar la búsqueda exacta de cliques \parencite{pardalos1992}.

    \item \textbf{Procedimiento algorítmico central:} La resolución se basó en el método exacto de Ramificación y Acotamiento (\textit{Branch \& Bound}) \parencite{carraghan1990, pardalos1992}, ejecutado internamente por HiGHS con planos de corte (Gomory, MIR) y heurísticas de factibilidad.

    \item \textbf{Cálculo de la relajación:} En cada nodo del árbol B\&B, el solver resuelve la relajación lineal continua del subproblema entero \parencite{nemhauser1988}, obteniendo una cota superior $z^*_{LP}$ del óptimo entero del subproblema.

    \item \textbf{Proceso de ramificación:} Cuando la solución de la relajación contiene una variable fraccionaria, el solver selecciona una variable $x_k$ con valor no entero y crea dos subproblemas mutuamente excluyentes ($x_k = 0$ y $x_k = 1$) \parencite{carraghan1990, nemhauser1988}. Los subproblemas se podan por infactibilidad, por cota (cuando $z^*_{LP}$ no mejora la incumbente) o por encontrar una solución entera factible.
\end{enumerate}

\subsection{Dificultades experimentales}\label{subsec:dificultades}

Durante la implementación y evaluación del modelo se enfrentaron tres dificultades principales, junto con la forma en que se abordaron:

\begin{enumerate}
    \item \textbf{Debilidad de la relajación LP.} La formulación ILP estándar (Ecuaciones \eqref{eq:obj}--\eqref{eq:bin}) genera una relajación lineal cuya cota es muy floja respecto al óptimo entero (brecha del 87{,}1\,\% en \textit{keller4}, ver Tabla~\ref{tab:relajacion}). Esto se abordó delegando la exploración del árbol Branch \& Bound al solver HiGHS, que combina planos de corte (Gomory, MIR) y heurísticas de factibilidad para cerrar la brecha mucho antes que un esquema B\&B puro.

    \item \textbf{Falta de convergencia de GLPK.} En las pruebas iniciales, GLPK no logró probar optimalidad sobre \textit{keller4} dentro de un límite de 300\,s, terminando con estado \textit{feasible}. La dificultad se manejó contrastando el resultado de HiGHS contra el óptimo conocido del benchmark DIMACS ($\omega = 11$) para certificar la validez de la respuesta, y reservando GLPK únicamente como punto de comparación de desempeño entre solvers.

    \item \textbf{Validación post-hoc del clique.} Para descartar errores en la lectura de la solución entera, se implementó una verificación del subgrafo inducido por los vértices seleccionados, comprobando que contiene exactamente $\binom{\omega}{2}$ aristas, todas presentes en el grafo original $G$.
\end{enumerate}

\begin{algorithm}
\caption{Ramificación y Acotamiento (Branch \& Bound) para el MCP}
\label{alg:branch_bound_mcp}
\algrenewcommand{\algorithmicindent}{1.0em}%
\algrenewcommand{\algorithmiccomment}[1]{\hfill$\triangleright$ \textit{#1}}%
\begin{algorithmic}[1]
\Require Grafo $G=(V,E)$, formulación ILP del MCP
\Ensure Clique máximo $x^*$ con cardinalidad $\check{z}$
\State $L \gets \{ P_0 \}$ \Comment{Lista de subproblemas activos}
\State $\check{z} \gets 0$ \Comment{Mejor solución entera}
\State $x^* \gets \emptyset$
\While{$L \neq \emptyset$}
\State Seleccionar y remover un subproblema $P$ de $L$
\State Resolver relajación LP de $P$: obtener $x'$, $z'$
\If{$P$ es infactible}
\State \textbf{descartar} $P$ \Comment{Poda por infactibilidad}
\ElsIf{$z' \leq \check{z}$}
\State \textbf{descartar} $P$ \Comment{Poda por cota}
\ElsIf{$x'$ es entera}
\State $\check{z} \gets z'$, $x^* \gets x'$ \Comment{Actualizar incumbente}
\Else
\State Elegir variable fraccionaria $x_k$ de $x'$
\State $L \gets L \cup \{P|_{x_k=0},\ P|_{x_k=1}\}$ \Comment{Ramificar}
\EndIf
\EndWhile
\State \Return $\check{z}, x^*$
\end{algorithmic}
\end{algorithm}

\FloatBarrier

\subsection{Diagrama de Flujo del Proceso}

La Figura \ref{fig:flowchart} presenta el diagrama de flujo del procedimiento implementado en el notebook experimental. Se utiliza Pyomo como framework de modelado algebraico y HiGHS como solver MIP, el cual resuelve internamente mediante Branch \& Bound con relajación LP, cortes y heurísticas.

\begin{figure}[ht]
\centering
\shorthandoff{>}%
\begin{tikzpicture}[
    scale=0.85, transform shape,
    node distance=1.1cm and 1.8cm,
    >={Stealth[length=3mm]},
    startstop/.style={rectangle, rounded corners=8pt, minimum width=3.2cm, minimum height=0.8cm, align=center, draw=black, fill=red!20, font=\small},
    process/.style={rectangle, minimum width=3.8cm, minimum height=0.8cm, align=center, draw=black, fill=blue!15, font=\small},
    solver/.style={rectangle, minimum width=3.8cm, minimum height=0.8cm, align=center, draw=black, fill=orange!20, font=\small},
    decision/.style={diamond, aspect=2.5, minimum width=1cm, align=center, draw=black, fill=yellow!20, font=\small, inner sep=1pt},
    io/.style={trapezium, trapezium left angle=70, trapezium right angle=110, minimum width=3cm, minimum height=0.8cm, align=center, draw=black, fill=green!15, font=\small},
    arrow/.style={thick, ->},
    dashedarrow/.style={thick, ->, dashed}
]

% --- Flujo principal (columna izquierda) ---
\node (start) [startstop] {Inicio};
\node (input) [io, below=of start] {Cargar instancia DIMACS\\$G=(V,E)$ (keller4)};
\node (complement) [process, below=of input] {Calcular grafo complemento\\$\bar{E} = \{(i,j) : (i,j) \notin E\}$};
\node (model) [process, below=of complement] {Formular modelo ILP (Pyomo)\\$\max \sum x_i$, s.a.\ $x_i+x_j \leq 1$};

% --- Resolucion ---
\node (solve) [solver, below=of model] {Resolver con HiGHS\\(Branch \& Bound interno)};
\node (status) [decision, below=of solve] {¿Óptimo?};

% --- Rama exito ---
\node (extract) [process, below=of status] {Extraer solución $x^*$\\$\omega = \sum x_i^*$};
\node (validate) [process, below=of extract] {Validar clique:\\$|E(S)| = \binom{\omega}{2}$?};

% --- Salida ---
\node (output) [io, below=of validate] {Generar visualizaciones:\\grafo con clique, grados};
\node (end) [startstop, below=of output] {Fin};

% --- Rama fallo (derecha cercana) ---
\node (fail) [startstop, right=2cm of status] {Error / Infactible};

% --- Relajacion LP (derecha lejana) ---
\node (lp) [solver, right=2.5cm of model] {Resolver relajación LP\\$x_i \in [0,1]$};
\node (gap) [process, below=of lp] {Calcular brecha de\\integralidad};

% === FLECHAS ===
\draw [arrow] (start) -- (input);
\draw [arrow] (input) -- (complement);
\draw [arrow] (complement) -- (model);
\draw [arrow] (model) -- (solve);
\draw [arrow] (solve) -- (status);
\draw [arrow] (status) -- node[right, font=\small] {Sí} (extract);
\draw [arrow] (status) -- node[above, font=\small] {No} (fail);
\draw [arrow] (extract) -- (validate);
\draw [arrow] (validate) -- (output);
\draw [arrow] (output) -- (end);

% Rama LP (desde modelo hacia la derecha)
\draw [arrow] (model.east) -- (lp.west);
\draw [arrow] (lp) -- (gap);
\draw [dashedarrow] (gap.south) |- ([yshift=-0.3cm]output.east) -- (output.east);

% Nota interna del solver
\node[below=0.1cm of solve, font=\tiny\itshape, text=gray] {LP relax + cortes + heurísticas};

\end{tikzpicture}%
\shorthandon{>}%
\caption{Diagrama de flujo de la implementación con Pyomo y HiGHS para el MCP.}
\label{fig:flowchart}
\end{figure}